% SPÖ Infografik-Poster: Wirtschaftliche Ordnung für Österreich
% Format: Portrait (A3/A2 druckfähig)
% Version: 1.0 | 2026-02-06
% Kompilierung: pdflatex POSTER_inflation_ordnung_2026-02-06.tex
% Basiert auf: NotebookLM Infografik Score 9.0
\documentclass[border=0pt]{standalone}
\usepackage[utf8]{inputenc}
\usepackage[T1]{fontenc}
\usepackage[ngerman]{babel}
\usepackage{tikz}
\usepackage{graphicx}
\usepackage{amssymb}
\usepackage{eurosym}
\usepackage{helvet}
\renewcommand{\familydefault}{\sfdefault}

\usetikzlibrary{shapes.geometric, positioning, calc, decorations.pathmorphing, patterns}

% === SPÖ FARBEN (Design System) ===
\definecolor{spored}{HTML}{E3000F}
\definecolor{spodarkred}{HTML}{B8000C}
\definecolor{spoblack}{HTML}{1A1A1A}
\definecolor{spodarkgray}{HTML}{333333}
\definecolor{spomedgray}{HTML}{666666}
\definecolor{spolightgray}{HTML}{F5F5F5}
\definecolor{spowhite}{HTML}{FFFFFF}
\definecolor{spogreen}{HTML}{2E8B57}
\definecolor{spolinegray}{HTML}{DDDDDD}

\begin{document}

\begin{tikzpicture}[
    every node/.style={inner sep=0, outer sep=0},
]

% === DIMENSIONS (portrait, ~A3 proportions) ===
\def\posterw{28}   % width in cm
\def\posterh{50}   % height in cm

% === BACKGROUND ===
\fill[spowhite] (0,0) rectangle (\posterw, \posterh);

% === SECTION 1: HEADER / TITLE ===
\node[anchor=north, text width=24cm, align=center] at (\posterw/2, \posterh-1.2) {
    {\fontsize{28}{34}\selectfont\bfseries\color{spoblack} Ordnen statt Spalten:}\\[8pt]
    {\fontsize{26}{32}\selectfont\bfseries\color{spoblack} Wirtschaftliche Ordnung f\"ur \"Osterreich}
};

% === SECTION 2: ALT vs NEU (Chaos vs Ordnung) ===
\def\sectwoy{42}

% Left box: ALT
\fill[spolightgray, rounded corners=6pt] (1.5, \sectwoy-6) rectangle (13.2, \sectwoy+1.2);

% Chaos scribbles (stylized tangled lines)
\begin{scope}[shift={(7.35, \sectwoy+0.1)}, scale=0.55]
    \draw[spodarkgray, line width=1.4pt, decorate, decoration={random steps, segment length=4pt, amplitude=2.5pt}]
        (-4,0.3) -- (-2,1.8) -- (0,-0.3) -- (2,1.2) -- (4,0.5);
    \draw[spodarkgray, line width=1.4pt, decorate, decoration={random steps, segment length=4pt, amplitude=2.5pt}]
        (-3.5,0.8) -- (-1,-1) -- (1.5,1) -- (3.5,-0.5);
    \draw[spodarkgray, line width=1.4pt, decorate, decoration={random steps, segment length=4pt, amplitude=2.5pt}]
        (-3,-0.8) -- (0,1.5) -- (3,-1);
    \draw[spodarkgray, line width=1.4pt, decorate, decoration={random steps, segment length=4pt, amplitude=2.5pt}]
        (-2.5,1.2) -- (0.5,-0.5) -- (2.5,0.8);
\end{scope}

\node[anchor=north, text width=10cm, align=center] at (7.35, \sectwoy-2.2) {
    {\fontsize{18}{22}\selectfont\bfseries\color{spored} ALT}\\[6pt]
    {\fontsize{11}{14}\selectfont\color{spodarkgray} Wirtschaftliches Chaos \&\\Populistische Versprechen}
};

% Vertical divider
\draw[spolinegray, line width=1pt] (14, \sectwoy-6) -- (14, \sectwoy+1.2);

% Right box: NEU
\fill[spolightgray, rounded corners=6pt] (14.8, \sectwoy-6) rectangle (26.5, \sectwoy+1.2);

% Order lines (stylized EKG/signal — clean)
\begin{scope}[shift={(20.65, \sectwoy+0.1)}, scale=0.55]
    \draw[spored, line width=2.5pt, line cap=round, line join=round]
        (-5,0) -- (-3,0) -- (-2.2,2) -- (-1.5,-2) -- (-0.8,2) -- (-0.1,-2) -- (0.6,0) -- (5,0);
\end{scope}

\node[anchor=north, text width=10cm, align=center] at (20.65, \sectwoy-2.2) {
    {\fontsize{18}{22}\selectfont\bfseries\color{spored} NEU}\\[6pt]
    {\fontsize{11}{14}\selectfont\color{spodarkgray} Stabile Ordnung \&\\Konkrete Massnahmen}
};

% === SECTION 3: KEY METRICS (3 columns) ===
\def\sectthreey{34}

% --- LEFT: Inflation sank ---
\fill[spowhite, rounded corners=6pt, draw=spolinegray, line width=0.6pt]
    (1.5, \sectthreey-6) rectangle (9.5, \sectthreey+1.5);

% Downward trend line
\begin{scope}[shift={(3.2, \sectthreey-1.3)}, scale=1]
    \draw[spolinegray, line width=0.5pt] (0,0) -- (0,3.8);
    \draw[spolinegray, line width=0.5pt] (0,0) -- (5,0);
    \draw[spored, line width=3pt, line cap=round, line join=round]
        (0.2,3.3) -- (1.2,2.8) -- (2.3,1.8) -- (3.3,0.8) -- (4.5,0.5);
    % Percentage labels
    \node[anchor=east, font=\fontsize{18}{22}\selectfont\bfseries, color=spoblack] at (-0.2, 3.3) {11\,\%};
    \node[anchor=west, font=\fontsize{18}{22}\selectfont\bfseries, color=spored] at (4.6, 0.5) {2\,\%};
\end{scope}

\node[anchor=north, text width=7cm, align=center, font=\fontsize{10}{13}\selectfont\color{spomedgray}]
    at (5.5, \sectthreey-4.7) {Inflation sank von 11\,\% auf 2\,\%.};

% --- CENTER: Starburst ---
\begin{scope}[shift={(\posterw/2, \sectthreey-2)}]
    % Starburst rays
    \foreach \i in {0,12,...,348} {
        \fill[spored] (0,0) -- (\i:3.2) -- (\i+6:2.4) -- cycle;
    }
    \fill[spored] (0,0) circle (2.4);

    % Text inside starburst
    \node[text width=4.2cm, align=center] at (0, 0.5) {
        {\fontsize{14}{17}\selectfont\bfseries\color{spowhite} 0,7 Prozent}\\[3pt]
        {\fontsize{14}{17}\selectfont\bfseries\color{spowhite} Preisr\"uckgang}\\[3pt]
        {\fontsize{14}{17}\selectfont\bfseries\color{spowhite} im J\"anner 2026}
    };
    \node[text width=4.2cm, align=center] at (0, -1.3) {
        {\fontsize{8}{10}\selectfont\color{spowhite!85!spored} Der st\"arkste monatliche}\\[1pt]
        {\fontsize{8}{10}\selectfont\color{spowhite!85!spored} R\"uckgang seit J\"anner 2019}
    };
\end{scope}

% --- RIGHT: Energiepreise ---
\fill[spowhite, rounded corners=6pt, draw=spolinegray, line width=0.6pt]
    (18.5, \sectthreey-6) rectangle (26.5, \sectthreey+1.5);

% Building + plug icon
\begin{scope}[shift={(22.5, \sectthreey+0.2)}, scale=0.45]
    % Building outline
    \draw[spodarkgray, line width=2pt] (-2,-2) rectangle (2,2);
    % Windows
    \foreach \wx in {-1,0.4} {
        \foreach \wy in {-0.8,0.6} {
            \fill[spolightgray] (\wx, \wy) rectangle (\wx+0.8, \wy+0.8);
            \draw[spodarkgray, line width=0.8pt] (\wx, \wy) rectangle (\wx+0.8, \wy+0.8);
        }
    }
    % Plug icon to the right
    \draw[spored, line width=2.5pt] (4,0) circle (0.8);
    \draw[spored, line width=2pt] (4,0.4) -- (4,0.9);
    \draw[spored, line width=2pt] (4,-0.4) -- (4,-0.9);
    \draw[spored, line width=1.5pt] (3.2,0) -- (2.5,0);
\end{scope}

\node[anchor=north, text width=7cm, align=center] at (22.5, \sectthreey-1.5) {
    {\fontsize{12}{15}\selectfont\bfseries\color{spoblack} Energiepreise um}\\[3pt]
    {\fontsize{18}{22}\selectfont\bfseries\color{spored} 4,9\,\% gesunken}\\[6pt]
    {\fontsize{10}{13}\selectfont\color{spodarkgray} Rund 300 Euro j\"ahrliche}\\[1pt]
    {\fontsize{10}{13}\selectfont\color{spodarkgray} Ersparnis pro Haushalt}\\[1pt]
    {\fontsize{10}{13}\selectfont\color{spodarkgray} bei Heizen und Tanken.}
};

% === SECTION 4: EHRLICHKEITS-LINIE (Quote) ===
\def\secfoury{26.5}

\fill[spolightgray, rounded corners=8pt] (2, \secfoury-3.8) rectangle (26, \secfoury+0.5);

\node[text width=21cm, align=center] at (\posterw/2, \secfoury-0.5) {
    {\fontsize{15}{19}\selectfont\bfseries\itshape\color{spoblack}
    \glqq Gegen die Teuerung gewinnt man nie ganz.}\\[4pt]
    {\fontsize{15}{19}\selectfont\bfseries\itshape\color{spoblack}
    Aber Schritt f\"ur Schritt -- mit Ordnung.\grqq}
};

\node[text width=21cm, align=center] at (\posterw/2, \secfoury-2.8) {
    {\fontsize{10}{13}\selectfont\color{spomedgray}
    R\"uckkehr zur Stabilit\"at durch klare Regeln statt leerer Versprechen.}
};

% === SECTION 5: DREI SCHRITTE ZUR ENTLASTUNG ===
\def\secfivey{21}

% Section divider with title
\draw[spolinegray, line width=1pt] (3, \secfivey+1.5) -- (9, \secfivey+1.5);
\draw[spolinegray, line width=1pt] (19, \secfivey+1.5) -- (25, \secfivey+1.5);
\node[fill=spowhite, inner sep=8pt, font=\fontsize{15}{18}\selectfont\bfseries\color{spoblack}]
    at (\posterw/2, \secfivey+1.5) {Drei Schritte zur Entlastung};

% Step 1
\fill[spolightgray, rounded corners=4pt] (3, \secfivey-0.6) rectangle (25, \secfivey+0.6);
\fill[spored!15!white, rounded corners=12pt] (3.8, \secfivey-0.35) circle (0.4);
\node[font=\fontsize{14}{16}\selectfont\bfseries, color=spored] at (3.8, \secfivey) {$\checkmark$};
\node[anchor=west, font=\fontsize{14}{17}\selectfont\bfseries\color{spoblack}] at (5, \secfivey) {
    1,4 Mrd.\ \euro{} MwSt-Senkung
};

% Step 2
\fill[spolightgray, rounded corners=4pt] (3, \secfivey-2.4) rectangle (25, \secfivey-1.2);
\fill[spored!15!white, rounded corners=12pt] (3.8, \secfivey-1.75) circle (0.4);
\node[font=\fontsize{14}{16}\selectfont\bfseries, color=spored] at (3.8, \secfivey-1.8) {$\checkmark$};
\node[anchor=west, font=\fontsize{14}{17}\selectfont\bfseries\color{spoblack}] at (5, \secfivey-1.8) {
    Mietpreisbremse f\"ur 300.000 Haushalte
};

% Step 3
\fill[spolightgray, rounded corners=4pt] (3, \secfivey-4.2) rectangle (25, \secfivey-3);
\fill[spored!15!white, rounded corners=12pt] (3.8, \secfivey-3.55) circle (0.4);
\node[font=\fontsize{14}{16}\selectfont\bfseries, color=spored] at (3.8, \secfivey-3.6) {$\checkmark$};
\node[anchor=west, font=\fontsize{14}{17}\selectfont\bfseries\color{spoblack}] at (5, \secfivey-3.6) {
    Anti-Shrinkflation-Gesetz
};

% === SECTION 6: VERGLEICH Österreich vs. Ungarn ===
\def\secsixy{13.5}

% Section divider with title
\draw[spolinegray, line width=1pt] (3, \secsixy+2) -- (7.5, \secsixy+2);
\draw[spolinegray, line width=1pt] (20.5, \secsixy+2) -- (25, \secsixy+2);
\node[fill=spowhite, inner sep=8pt, font=\fontsize{15}{18}\selectfont\bfseries\color{spoblack}]
    at (\posterw/2, \secsixy+2) {Aktive Markt-Ordnung vs.\ Populismus};

% Left card: SPÖ-Modell (Österreich)
\fill[spolightgray, rounded corners=8pt] (2.5, \secsixy-4.5) rectangle (13.5, \secsixy+0.8);

% Green checkmark badge
\fill[spogreen] (4.8, \secsixy+0.1) circle (0.55);
\node[font=\fontsize{14}{16}\selectfont\bfseries\color{spowhite}] at (4.8, \secsixy+0.1) {$\checkmark$};
\node[anchor=west, font=\fontsize{12}{14}\selectfont\bfseries\color{spoblack}] at (5.7, \secsixy+0.1) {
    SP\"O-Modell (\"Osterreich)
};

\node[anchor=north, text width=9cm, align=center] at (8, \secsixy-1) {
    {\fontsize{24}{28}\selectfont\bfseries\color{spoblack} 2,0\,\%}\\[2pt]
    {\fontsize{14}{17}\selectfont\bfseries\color{spoblack} Inflationsrate}\\[6pt]
    {\fontsize{11}{14}\selectfont\color{spomedgray} Aktive Markt-Ordnung}
};

% Right card: Populismus (Ungarn)
\fill[spolightgray, rounded corners=8pt] (14.5, \secsixy-4.5) rectangle (25.5, \secsixy+0.8);

% Red X badge
\fill[spored] (16.8, \secsixy+0.1) circle (0.55);
\node[font=\fontsize{14}{16}\selectfont\bfseries\color{spowhite}] at (16.8, \secsixy+0.1) {$\times$};
\node[anchor=west, font=\fontsize{12}{14}\selectfont\bfseries\color{spoblack}] at (17.7, \secsixy+0.1) {
    Populismus (Ungarn)
};

\node[anchor=north, text width=9cm, align=center] at (20, \secsixy-1) {
    {\fontsize{24}{28}\selectfont\bfseries\color{spoblack} 4,0\,\%}\\[2pt]
    {\fontsize{14}{17}\selectfont\bfseries\color{spoblack} Inflationsrate}\\[6pt]
    {\fontsize{11}{14}\selectfont\color{spomedgray} \"Okonomische Instabilit\"at}
};

% === SECTION 7: HASHTAG CTA ===
\def\secseveny{6.5}

\node[anchor=north, text width=24cm, align=center] at (\posterw/2, \secseveny+0.5) {
    {\fontsize{38}{44}\selectfont\bfseries\color{spored} \#OrdnenStattSpalten}
};

\node[anchor=north, text width=24cm, align=center] at (\posterw/2, \secseveny-2.5) {
    {\fontsize{14}{18}\selectfont\color{spoblack} Wer ordnet, handelt. Wer handelt, entlastet.}
};

% === SECTION 8: FOOTER / QUELLEN ===
\draw[spolinegray, line width=0.4pt] (3, 2) -- (25, 2);

\node[anchor=north, text width=22cm, align=center, font=\fontsize{8}{10}\selectfont\color{spomedgray}]
    at (\posterw/2, 1.7) {
    Quellen: Statistik Austria, WIFO, Eurostat, Wirtschaftsforschungsinstitute. Stand: Februar 2026.
};

\end{tikzpicture}

\end{document}
